\chapter{The Canonical Ensemble}

\section{Motivation: How Can Temperature Be the Same if Energy Differs?}

Consider a \emph{system} connected to a large \emph{thermal reservoir} at fixed temperature $T$.
The combined object — reservoir plus system — is an isolated total system, so the
microcanonical ensemble applies to it.  The total extensive quantities satisfy
\begin{equation}
    U_{\text{tot}} = U_R + U_S, \qquad
    V_{\text{tot}} = V_R + V_S, \qquad
    N_{\text{tot}} = N_R + N_S, \qquad
    S_{\text{tot}} = S_R + S_S.
\end{equation}

\begin{center}[DIAGRAM: reservoir at temperature $T$ on the left exchanging heat with a small system on the right; arrows labeled ``heat'']\end{center}

The key question is: what is the probability that we find the \emph{system} in a specific
microstate $j$ that carries energy $E_j$?  Note carefully that we are asking for the probability of
a particular microstate of the system, \emph{not} merely for the probability that the system has
energy $E_j$.

\subsection{Deriving the Boltzmann Factor}

The key idea is that even though the system's energy fluctuates, the combined system is
isolated with fixed total energy $U$.  By the fundamental postulate of statistical mechanics,
every microstate of the \emph{total} system is equally probable.  Therefore, the probability
of finding the system in microstate $j$ (with energy $E_j$) is proportional to the number of
microstates available to the \emph{reservoir} when it has energy $U - E_j$:
\begin{equation}
    P_j(E_j) \;\propto\; \Gamma_S(j)\;\times\;\Gamma_R(U - E_j),
\end{equation}
where $\Gamma_S(j) = 1$ because we have fixed the system in a single microstate $j$, and
$\Gamma_R(U-E_j)$ counts the number of ways the reservoir can have energy $U - E_j$.  Thus
\begin{equation}
    P_j(E_j) \;\propto\; \Gamma_R(U - E_j) \;=\; e^{S_R(U - E_j)/k_B}.
\end{equation}

Because the reservoir is large, $E_j \ll U$, and we may Taylor-expand $S_R(U - E_j)$ about
$U$:
\begin{equation}
    S_R(U - E_j) \;\approx\; S_R(U) - \left.\pdv{S_R}{U}\right|_V E_j
    \;=\; S_R(U) - \frac{E_j}{T},
\end{equation}
where we used the thermodynamic identity $(\partial S_R/\partial U)_V = 1/T$.  The constant
$S_R(U)$ is the same for every microstate $j$ and drops into an overall normalisation factor.
We are left with the celebrated \textbf{Boltzmann factor}:

\begin{formula}[Boltzmann factor]
\begin{equation}
    P_j(E_j) \;\propto\; e^{-E_j / k_B T}.
    \label{eq:boltzmann_factor}
\end{equation}
\end{formula}

The validity of the first-order Taylor expansion requires that higher-order terms be
negligible, i.e.\ $E_j \ll U_R$.  A rough check shows that the next term is proportional to
$E_j\,\partial^2 S/\partial U^2$, which is of order $E_j/(k_B T^2)\cdot k_B = E_j/T^2$,
so the correction is suppressed by the factor $E_j/(k_B T) \ll 1$ for states that are
appreciably occupied.

\subsection{The Partition Function and Normalisation}

Requiring that the probabilities sum to one determines the normalisation.  Writing
$\beta \equiv 1/(k_B T)$ for convenience,
\begin{equation}
    P_j(E_j) = \frac{e^{-\beta E_j}}{Z}, \qquad
    Z \equiv \sum_j e^{-\beta E_j}.
\end{equation}

\begin{formula}[Partition function]
\begin{equation}
    Z(T,V,N) \;=\; \sum_j e^{-\beta E_j},
    \label{eq:partition_function}
\end{equation}
where the sum runs over \emph{all microstates} $j$ of the system (not just over distinct
energies).
\end{formula}

The partition function $Z$ is a new state function depending on $T$, $V$, and $N$.  It is
sometimes called the \emph{canonical partition function} to distinguish it from other ensemble
partition functions introduced later.

\section{The Canonical Ensemble}

\begin{definition}[Canonical ensemble]
The canonical ensemble is the ensemble of microstates of a closed system in equilibrium with a
heat reservoir at temperature $T$.  The symbols `microstates' refer to all microstates at all
energies (unlike the microcanonical ensemble, which is restricted to a single energy shell),
and the probability of each microstate $j$ is
\begin{equation}
    P_j = \frac{1}{Z}\,e^{-\beta E_j}.
\end{equation}
\end{definition}

\section{Calculating the Partition Function}

In practice the sum in \eqref{eq:partition_function} is rarely performed by listing individual
quantum numbers.  It is more natural to group states by energy.  If $g(E_j)$ denotes the
degeneracy (number of microstates with energy $E_j$), the sum over microstates becomes a sum
over distinct energies:
\begin{equation}
    Z \;=\; \sum_j e^{-\beta E_j}
      \;=\; \sum_{E_i} g(E_i)\,e^{-\beta E_i}.
\end{equation}

In the limit of a continuous energy spectrum the degeneracy becomes a density of states
$dN/dE$, and the sum becomes an integral:
\begin{equation}
    Z \;=\; \int dE\;\frac{dN}{dE}\;e^{-\beta E}.
\end{equation}

For systems with a classical (semi-classical) Hamiltonian, the density of states is obtained by
integrating over phase space.  For a single particle in three dimensions,
\begin{equation}
    Z \;=\; \int \prod_{i=1}^{3} \frac{d^3x_i\,d^3p_i}{(2\pi\hbar)^3}\;e^{-\beta E(\vec{x},\vec{p})}.
\end{equation}

\section{Thermodynamics from the Partition Function (The Thermodynamics Recipe)}

The canonical ensemble provides a systematic recipe for extracting all thermodynamic quantities
from $Z$.

\subsection{Internal Energy}

The internal energy is the ensemble average of the energy:
\begin{equation}
    U = \langle E \rangle
      = \frac{\sum_j E_j\,P_j}{\sum_j P_j}
      = \frac{\sum_j E_j\,e^{-\beta E_j}}{\sum_j e^{-\beta E_j}}.
\end{equation}
Observing that $E_j\,e^{-\beta E_j} = -\partial e^{-\beta E_j}/\partial\beta$, the numerator
is $-\partial Z/\partial\beta$, so
\begin{equation}
    \langle E \rangle = -\frac{1}{Z}\pdv{Z}{\beta} = -\pdv{}{\beta}\log Z.
\end{equation}

\begin{formula}[Internal energy from partition function]
\begin{equation}
    U = \langle E \rangle = -\pdv{\log Z}{\beta}.
    \label{eq:U_from_Z}
\end{equation}
\end{formula}

It is worth noting that $\partial/\partial\beta = (\partial T/\partial\beta)\,\partial/\partial T = -k_B T^2\,\partial/\partial T$, so $\partial \log Z/\partial\beta = -k_B T^2\,\partial\log Z/\partial T$.

\subsection{Entropy}

Starting from the Gibbs entropy $S = -k_B \sum_j P_j \log P_j$ with $P_j = e^{-\beta E_j}/Z$,
we substitute $\log P_j = -\beta E_j - \log Z$:
\begin{align}
    S &= -k_B \sum_j \frac{e^{-\beta E_j}}{Z}\bigl(-\beta E_j - \log Z\bigr) \notag\\
      &= k_B\beta \sum_j \frac{E_j\,e^{-\beta E_j}}{Z} + k_B\log Z \sum_j \frac{e^{-\beta E_j}}{Z} \notag\\
      &= k_B\beta\,U + k_B\log Z.
\end{align}

\begin{formula}[Entropy from partition function]
\begin{equation}
    S = \frac{U}{T} + k_B \log Z.
    \label{eq:S_from_Z}
\end{equation}
\end{formula}

\subsection{Helmholtz Free Energy}

From \eqref{eq:S_from_Z} we immediately see that $U - TS = -k_B T\log Z$.  Since $F \equiv U - TS$ is the Helmholtz free energy,

\begin{formula}[Helmholtz free energy]
\begin{equation}
    F = -k_B T\log Z.
    \label{eq:F_from_Z}
\end{equation}
\end{formula}

All other thermodynamic quantities follow from $F$.  The thermodynamic relations $dF = -SdT - PdV$ give:

\begin{formula}[Equation of state]
\begin{equation}
    S = -\left.\pdv{F}{T}\right|_V = \frac{\partial(k_B T\log Z)}{\partial T}\bigg|_V, \qquad
    P = -\left.\pdv{F}{V}\right|_T.
\end{equation}
\end{formula}

\subsection{Heat Capacity}

The heat capacity at constant volume is
\begin{align}
    C_V &= \left.\pdv{U}{T}\right|_V
         = \pdv{}{T}\left[-\pdv{\log Z}{\beta}\right]
         = \pdv{}{T}\!\left[\frac{1}{k_B T^2}\pdv{\log Z}{\beta}\cdot k_B T^2\right]\notag\\
         &= -k_B\beta^2\,\pdv[2]{\log Z}{\beta}\bigg|_V
\end{align}
or equivalently $C_V = k_B\beta^2\,\partial^2\log Z/\partial\beta^2$.

\section{Fluctuations in Energy}

The heat capacity also controls the size of energy fluctuations.  Recall that
$\text{Var}(E) = \langle E^2\rangle - \langle E\rangle^2$.  Differentiating $U = -\partial\log Z/\partial\beta$ a second time gives
\begin{equation}
    \pdv[2]{\log Z}{\beta} = \frac{\partial^2 Z/\partial\beta^2}{Z} - \left(\frac{\partial Z/\partial\beta}{Z}\right)^2
    = \langle E^2\rangle - \langle E\rangle^2 = \text{Var}(E).
\end{equation}
Comparing with the expression for $C_V$:

\begin{formula}[Energy fluctuations]
\begin{equation}
    \text{Var}(E) \;=\; k_B T^2\, C_V.
    \label{eq:var_E}
\end{equation}
\end{formula}

This is a striking result: the equilibrium fluctuations of the energy are directly related to the
macroscopic response function $C_V$.  It is an example of the fluctuation-response theorem.

\subsection*{Summary of the Thermodynamics Recipe}

\begin{theorem}[Canonical ensemble thermodynamics]
Given the partition function $Z(T,V,N)$:
\begin{alignat}{3}
    &(1)\quad U = -\pdv{\log Z}{\beta}, \qquad
    &&(3)\quad S = \frac{U}{T} + k_B\log Z, \qquad
    &&(5)\quad F = -k_B T\log Z,\\[4pt]
    &(2)\quad P = -\left.\pdv{F}{V}\right|_T, \qquad
    &&(4)\quad C_V = k_B\beta^2\pdv[2]{\log Z}{\beta}, \qquad
    &&(6)\quad \mathrm{Var}(U) = k_B T^2 C_V.
\end{alignat}
\end{theorem}

\subsection{Are the Fluctuations Appreciable?}

For an ideal gas, $C_V = \tfrac{3}{2}Nk_B$ and $U = \tfrac{3}{2}Nk_BT$, so the relative
fluctuation in energy is
\begin{equation}
    \frac{\sigma(U)}{U} = \frac{\sqrt{k_BT^2 C_V}}{U}
    = \frac{\sqrt{k_BT^2\cdot\tfrac{3}{2}Nk_B}}{\tfrac{3}{2}Nk_BT}
    = \sqrt{\frac{3}{2}}\cdot\frac{1}{\sqrt{N}}.
\end{equation}
Relative fluctuations are thus suppressed by $1/\sqrt{N}$ and are negligible in the
thermodynamic limit.  More generally, since $C_V \sim N$ and $U \sim N$, we always have
$\sigma(U)/U \sim 1/\sqrt{N}$.  This justifies treating the canonical and microcanonical
ensembles as equivalent for macroscopic systems.

\section{Example: The Two-State System}

\subsection{Partition Function and Internal Energy}

Consider a spin-$\tfrac{1}{2}$ particle in an external magnetic field $B$.  The two energy
levels are
\begin{equation}
    E_1 = -\mu_B B \quad\text{(spin aligned with field)}, \qquad
    E_2 = +\mu_B B \quad\text{(spin anti-aligned)}.
\end{equation}
The partition function is
\begin{equation}
    Z = e^{-\beta E_1} + e^{-\beta E_2}
      = e^{\beta\mu_B B} + e^{-\beta\mu_B B}
      = 2\cosh(\beta\mu_B B).
\end{equation}
The internal energy follows from \eqref{eq:U_from_Z}:
\begin{equation}
    U = -\pdv{\log Z}{\beta} = -\pdv{}{\beta}\log\bigl[2\cosh(\beta\mu_B B)\bigr]
      = -\mu_B B\tanh(\beta\mu_B B).
\end{equation}

\begin{center}[DIAGRAM: plot of $U$ vs.\ $T$; the curve rises from $-\mu_B B$ at low $T$ (all spins in ground state) to $0$ at high $T$ (equally likely in $E_1$ or $E_2$)]\end{center}

\subsection{Entropy}

From \eqref{eq:S_from_Z},
\begin{align}
    S &= \frac{U}{T} + k_B\log Z \notag\\
      &= k_B\log\bigl(2\cosh(\beta\mu_B B)\bigr) + k_BT\cdot\frac{1}{k_BT}\,\pdv{\log Z}{\beta}\cdot\beta
      \notag\\
      &= k_B\log\bigl(2\cosh(\beta\mu_B B)\bigr) - \frac{\mu_B B}{T}\tanh(\beta\mu_B B).
\end{align}

Checking the limits: as $T\to 0$, $\beta\to\infty$ and $\cosh(\beta\mu_B B)\to\tfrac{1}{2}e^{\beta\mu_B B}$,
so $\log(2\cosh)\to \beta\mu_B B$ and $\tanh\to 1$, giving $S\to k_B\cdot\beta\mu_B B - k_B\cdot\beta\mu_B B = 0$
(the system is in a unique ground state, consistent with the third law).  As $T\to\infty$,
$\tanh\to 0$ and $\cosh\to 1$, so $S\to k_B\log 2$ (two equally accessible states).

\subsection{Heat Capacity}

Differentiating $U$ with respect to $T$,
\begin{equation}
    C_V = \pdv{U}{T} = \pdv{}{T}\bigl[-\mu_B B\tanh(\beta\mu_B B)\bigr]
        = \frac{(\mu_B B)^2}{k_B T^2}\,\mathrm{sech}^2(\beta\mu_B B)
        = \frac{(\mu_B B)^2}{k_BT^2}\cdot\frac{1}{\cosh^2(\beta\mu_B B)}.
\end{equation}
This is the \emph{Schottky anomaly}: the heat capacity peaks at an intermediate temperature
$k_BT \sim \mu_B B$ (where the thermal energy is comparable to the level spacing) and falls
off exponentially at low $T$ and as $T^{-2}$ at high $T$.

\subsection{Energy Fluctuations}

The fluctuation in energy is $\sigma(U) = \sqrt{k_BT^2 C_V}$:
\begin{equation}
    \sigma(U) = \frac{\mu_B B}{\cosh(\beta\mu_B B)}\,\mu_B B \cdot \frac{1}{k_BT}.
\end{equation}
In the limits: $T\to 0 \Rightarrow \sigma(U)\to 0$ (system certain to be in $E_1$); $T\to\infty
\Rightarrow \sigma(U)\to\mu_B B$ (equally likely to be in $E_1$ or $E_2$).

\subsection{Individual State Probabilities}

We may also extract the probabilities $P_1$ and $P_2$ directly:
\begin{align}
    P_1 &= \frac{e^{\beta\mu_B B}}{Z} = \frac{e^{\beta\mu_B B}}{e^{\beta\mu_B B}+e^{-\beta\mu_B B}}
          = \frac{1}{1+e^{-2\beta\mu_B B}}, \\[4pt]
    P_2 &= \frac{1}{1+e^{2\beta\mu_B B}}.
\end{align}
As a consistency check, $P_1 + P_2 = 1$.

\chapter{Recitation: Carnot Engine and Surface Tension}

\section{Carnot Engine Revisited}

Consider two heat reservoirs: a hot reservoir at temperature $T_H$ and a cold reservoir at
temperature $T_C < T_H$.  A Carnot engine operates between them with a working substance
(total heat capacity $C_V$) that is initially at temperature $T_2$.

\begin{center}[DIAGRAM: left panel shows the $P$-$V$ Carnot cycle with isotherms and adiabats labeled, and arrows indicating $Q_{\text{in}}$, $Q_{\text{out}}$, and work; right panel shows the $T$-$S$ diagram for the same cycle, with processes labeled 1--4 and the note that the system does not return to the same temperature in this example]\end{center}

The entropy change of the hot reservoir as the working substance absorbs heat $\delta Q_{\text{in}}$ is $dS_H = -\delta Q_{\text{in}}/T$ (the reservoir loses entropy).  While the system
temperature rises from $T_2$ to $T_C$ in contact with $T_H$, the entropy change of the hot
reservoir is
\begin{equation}
    \Delta S_H = \int_{T_2}^{T_C} C_V\,\frac{dT_H}{T_H} = C_V\ln\frac{T_C}{T_2}
               = -C_V\ln\frac{T_2}{T_C}.
\end{equation}

The total net work done per cycle is
\begin{equation}
    W_{\text{net}} = Q_H - Q_C, \qquad \eta \equiv \frac{W_{\text{net}}}{Q_H}\;\text{(efficiency)},
\end{equation}
so
\begin{equation}
    \eta = \frac{Q_H - Q_C}{Q_H} = 1 - \frac{Q_C}{Q_H}.
\end{equation}
From the $T$-$S$ diagram (where $Q = T\Delta S$ for a reversible process), $Q_C = T_C\Delta S$
and $Q_H = T_H\Delta S$, giving the Carnot efficiency $\eta = 1 - T_C/T_H$.

Integrating over a full cycle where the hot reservoir drives the system from $T_2$ to $T_C$,
the total net work is
\begin{equation}
    W_{\text{net}} = C_V(T_2 - T_C) - C_V T_C\ln\frac{T_H}{T_C},
\end{equation}
where the first term is the ``amount of heat you'd like'' and the second is the ``waste heat.''

\section{Surface Tension Problem}

Water molecules at a surface experience less attraction from their neighbours than those in the
bulk, so the surface has higher energy.  Systems minimise their surface energy, which is why
droplets form spheres.

Consider a film of fluid supported by a rectangular frame with one moveable side of length
$\ell$.  The surface tension is
\begin{equation}
    \sigma = \sigma_0 - \alpha T,
\end{equation}
reflecting the empirical fact that surface tension decreases linearly with temperature (this
linear dependence is well-known for many liquids).  The force needed to extend the film by
$dx$ (a factor of 2 because the film has two surfaces) is $F = 2\sigma\ell$, so the work done
on the system is $\delta W = -2\sigma\ell\,dx$ (note the sign: stretching the film does work on
it).

\begin{center}[DIAGRAM: rectangular frame with width $\ell$ and one moveable side displaced by $x$; two surface layers indicated]\end{center}

The first law for this system is
\begin{equation}
    dU = \delta Q - 2\sigma\ell\,dx = T\,dS + 2\sigma\ell\,dx,
\end{equation}
where we set the work term with the sign appropriate to having $x$ as the extensive variable
(increasing $x$ increases the area, adding energy if $\sigma>0$).

To find $(\partial\sigma/\partial x)_T$ we use the Helmholtz free energy $F = U - TS$:
\begin{align}
    dF &= dU - T\,dS - S\,dT \notag\\
       &= 2\sigma\ell\,dx - S\,dT.
\end{align}
The partials must be the same:
\begin{equation}
    \left.\pdv{F}{x}\right|_T = 2\sigma\ell, \qquad
    \left.\pdv{F}{T}\right|_x = -S.
\end{equation}
Taking a mixed derivative and equating,
\begin{equation}
    \frac{\partial^2 F}{\partial T\,\partial x} = 2\ell\left.\pdv{\sigma}{T}\right|_x
    = -\left.\pdv{S}{x}\right|_T.
\end{equation}
Therefore
\begin{equation}
    \left.\pdv{U}{x}\right|_T = 2\ell\sigma + T\left.\pdv{S}{x}\right|_T
    = 2\ell\sigma - 2\ell T\left.\pdv{\sigma}{T}\right|_x
    = 2\ell(\sigma_0 - \alpha T) + 2\ell T\alpha = 2\ell\sigma_0.
\end{equation}
Thus $\Delta U = 2\sigma_0\ell\,\Delta x$.  Interestingly, the work done $2\sigma\ell\,\Delta x
= 2(\sigma_0-\alpha T)\ell\,\Delta x < \Delta U$, which means the system absorbs heat from the
reservoir upon stretching.  This is the \emph{opposite} of a rubber band, where stretching is
exothermic.

\chapter{Thermodynamics of the Ideal Gas}

\section{Computing the Ideal Gas Partition Function}

\subsection{Derivation}

For an ideal gas of $N$ distinguishable particles in volume $V$, the $N$-particle partition
function factorises into a product of single-particle partition functions because the Hamiltonian
is a sum of single-particle Hamiltonians:
\begin{equation}
    Z_N = \int\prod_{i=1}^{N}\frac{d^3x_i\,d^3p_i}{(2\pi\hbar)^3}\;
          e^{-\beta\sum_i p_i^2/2m}.
\end{equation}
The spatial integrals each contribute $V$, and the momentum integrals factorise:
\begin{align}
    Z_N &= \left(\frac{1}{2\pi\hbar}\right)^{3N} V^N
           \left(\int d^3p\;e^{-\beta p^2/2m}\right)^N.
\end{align}
Evaluating the momentum integral in spherical coordinates,
\begin{align}
    \int d^3p\;e^{-\beta p^2/2m}
    &= 4\pi\int_0^\infty p^2\,e^{-\beta p^2/2m}\,dp.
\end{align}
Substituting $u = \sqrt{\beta/2m}\,p$ so that $p^2 = (2m/\beta)u^2$ and $dp = \sqrt{2m/\beta}\,du$,
\begin{equation}
    = 4\pi\left(\frac{2m}{\beta}\right)^{3/2}\int_0^\infty u^2 e^{-u^2}\,du
    = 4\pi\left(\frac{2m}{\beta}\right)^{3/2}\cdot\frac{\sqrt{\pi}}{4}
    = \left(\frac{2\pi m}{\beta}\right)^{3/2}
    = (2\pi mk_BT)^{3/2}.
\end{equation}
Putting it together,
\begin{equation}
    Z_N = \left(\frac{1}{2\pi\hbar}\right)^{3N}V^N(2\pi mk_BT)^{3N/2}
        = \left(\frac{mk_BT}{2\pi\hbar^2}\right)^{3N/2}V^N
        = \frac{V^N}{\lambda_{\text{th}}^{3N}},
\end{equation}
where the \emph{thermal de Broglie wavelength} is
\begin{equation}
    \lambda_{\text{th}} = \frac{h}{\sqrt{2\pi mk_BT}}.
\end{equation}

\begin{formula}[Ideal gas partition function]
\begin{equation}
    Z_N = \left(\frac{V}{\lambda_{\text{th}}^3}\right)^N = (Z_1)^N,
    \label{eq:ideal_gas_Z}
\end{equation}
where $Z_1 = V/\lambda_{\text{th}}^3$ is the single-particle partition function.
\end{formula}

\subsection{Checking the Entropy: Gibbs Paradox}

From $F = -k_BT\log Z_N$ and $S = -(\partial F/\partial T)_V$,
\begin{align}
    S &= k_B\log Z_N + k_BT\frac{1}{Z_N}\pdv{Z_N}{T}\bigg|_V \notag\\
      &= k_BN\log V + k_BN\cdot\tfrac{3}{2}\log\!\left(\frac{mk_BT}{2\pi\hbar^2}\right)
         + k_BNT\cdot\frac{3}{2T} \notag\\
      &= k_BN\left[\log V + \frac{3}{2}\log\!\left(\frac{mk_BT}{2\pi\hbar^2}\right) + \frac{3}{2}\right].
\end{align}
This entropy is \emph{not extensive}: it depends on $V$ rather than on $V/N$, so doubling the
system at fixed density would double $S$ only if the $\log V$ term is replaced by $\log(V/N)$.
This is the \textbf{Gibbs paradox}.  The fix is to account for the indistinguishability of
identical particles by dividing the phase-space volume by $N!$, giving
\begin{equation}
    Z_N = \frac{1}{N!}(Z_1)^N, \qquad Z_N = \frac{1}{N!}\left(\frac{V}{\lambda_{\text{th}}^3}\right)^N.
\end{equation}
With this correction (using Stirling's approximation $\log N! \approx N\log N - N$), the entropy
becomes the \textbf{Sackur--Tetrode equation}:
\begin{formula}[Sackur--Tetrode equation]
\begin{equation}
    S = k_BN\left[\log\frac{V}{N\lambda_{\text{th}}^3} + \frac{3}{2}\log\!\left(\frac{V}{N\lambda_{\text{th}}^3}\right)\cdot\frac{0}{1} + \frac{5}{2}\right]
      = k_BN\left[\log\frac{V}{N\lambda_{\text{th}}^3} + \frac{5}{2}\right].
\end{equation}
\end{formula}
This is now properly extensive.

\subsection{Checking the Internal Energy}

Using \eqref{eq:U_from_Z} with $\log Z_N = N\log(V/\lambda_{\text{th}}^3) - \log N!$:
\begin{equation}
    U = -\pdv{\log Z_N}{\beta} = -N\pdv{}{\beta}\!\left[\log V - \tfrac{3}{2}\log\beta + \text{const}\right]
      = \frac{3N}{2\beta} = \frac{3}{2}Nk_BT.
\end{equation}
This recovers the equipartition result for a monatomic ideal gas with three translational degrees of freedom.

\section{Where Does the Equipartition Theorem Arise?}

The equipartition theorem is a direct consequence of the Gaussian structure of Boltzmann
integrals.  Suppose $\xi$ is any degree of freedom that enters the Hamiltonian quadratically:
$H = a\xi^2 + \cdots$.  The canonical integral over $\xi$ is
\begin{equation}
    Z = \int_{-\infty}^{\infty} d\xi\;e^{-\beta a\xi^2}\times\cdots
      = \sqrt{\frac{\pi}{\beta a}}\times\cdots
      = \sqrt{\frac{\pi k_BT}{a}}\times\cdots.
\end{equation}
The contribution to the internal energy from this degree of freedom is
\begin{equation}
    U_\xi = -\pdv{\log Z}{\beta} = -\pdv{}{\beta}\!\left[-\tfrac{1}{2}\log(\beta a) + \cdots\right]
           = \frac{1}{2\beta} = \frac{1}{2}k_BT.
\end{equation}
This is the equipartition contribution: each quadratic degree of freedom contributes $\tfrac{1}{2}k_BT$ to the average energy, regardless of the coefficient $a$.

\chapter{Probability Distributions from the Canonical Ensemble}

\section{Maxwell--Boltzmann Energy Distribution}

The canonical ensemble allows us to derive not just average quantities but the full probability
distributions for observables.  The probability density in phase space is
\begin{equation}
    \rho(\vec{x},\vec{p}) = \frac{d^3x\,d^3p}{(2\pi\hbar)^3}\cdot\frac{1}{Z}\,e^{-\beta E(\vec{x},\vec{p})},
\end{equation}
which is the joint probability density (probability per unit phase-space volume).

The energy distribution $P(E)$ is obtained by integrating over all phase-space points with
energy $E$, i.e.\ by inserting a delta function:
\begin{align}
    P(E) &= \int \frac{d^3x\,d^3p}{(2\pi\hbar)^3}\cdot\frac{1}{Z}\,e^{-\beta E(\vec{x},\vec{p})}
             \;\delta\!\left(E - \frac{p^2}{2m}\right) \notag\\
         &= \frac{V}{Z}\cdot\frac{1}{(2\pi\hbar)^3}\int d^3p\;e^{-\beta p^2/2m}\;\delta\!\left(E - \frac{p^2}{2m}\right).
\end{align}
Using $Z_1 = V/\lambda_{\text{th}}^3$ and evaluating the integral in spherical momentum coordinates
(the delta function localises to $p = \sqrt{2mE}$):
\begin{align}
    P(E) &= \frac{V}{Z_1}\cdot\frac{4\pi}{(2\pi\hbar)^3}\cdot p^2\,e^{-\beta p^2/2m}
             \cdot\frac{1}{|dE/dp|}_{p=\sqrt{2mE}} \notag\\
         &= \frac{V}{Z_1}\cdot\frac{4\pi}{(2\pi\hbar)^3}\cdot(2mE)\,e^{-\beta E}
             \cdot\frac{m}{\sqrt{2mE}}\cdot \frac{1}{2\sqrt{\beta E}} \cdot \text{(Jacobian)}.
\end{align}
After careful evaluation (substituting $u^2 = \beta E$ and using $\int_0^\infty u^2 e^{-u^2}du=\sqrt{\pi}/4$), the result is

\begin{formula}[Maxwell--Boltzmann energy distribution]
\begin{equation}
    P(E) = 2\pi\left(\frac{1}{\pi k_BT}\right)^{3/2}\sqrt{E}\;e^{-\beta E}.
    \label{eq:MB_energy}
\end{equation}
\end{formula}

This distribution peaks at $E = \tfrac{1}{2}k_BT$ and has mean $\langle E\rangle = \tfrac{3}{2}k_BT$.

\section{Maxwell--Boltzmann Speed Distribution}

The speed distribution $P(v)$ is obtained from the energy distribution by using $E = \tfrac{1}{2}mv^2$ and requiring that probability is conserved under the change of variable: $P(v)\,dv = P(E)\,dE$.

\begin{align}
    P(v) &= \int dE\;P(E)\;\delta\!\left(v - \sqrt{\frac{2E}{m}}\right) \notag\\
         &= 2\pi\left(\frac{1}{\pi k_BT}\right)^{3/2} \int_0^\infty dE\;\sqrt{E}\;e^{-\beta E}\;
            \delta\!\left(v - \sqrt{\frac{2E}{m}}\right).
\end{align}
The delta function sets $E = \tfrac{1}{2}mv^2$, contributing a Jacobian factor $|dv/dE| = 1/(mv)$:

\begin{formula}[Maxwell--Boltzmann speed distribution]
\begin{equation}
    P(v) = 4\pi\left(\frac{m}{2\pi k_BT}\right)^{3/2} v^2\,e^{-mv^2/2k_BT}.
    \label{eq:MB_speed}
\end{equation}
\end{formula}

\begin{center}[DIAGRAM: three-dimensional velocity space showing a sphere of radius $v_0$; the volume of the spherical shell $4\pi v^2 dv$ (the Boltzmann factor times the shell volume) gives $P(v)$; velocity components $v_x$, $v_y$, $v_z$ labeled]\end{center}

The $v^2$ prefactor reflects the growing phase-space volume of a spherical shell in velocity
space, while $e^{-mv^2/2k_BT}$ is the Boltzmann suppression at high speeds.

\chapter{Systems of Many Distinguishable and Indistinguishable Particles}

\section{Recall: Distinguishable Particles}

For an ideal gas of $N$ distinguishable atoms, we found $Z_N = (Z_1)^N$ because the
Hamiltonian is separable: $H = \sum_i H_i$, so the $N$-body phase-space integral factorises
into $N$ independent single-particle integrals.

\section{Generalisation: Many-Body Systems with Single-Particle Levels}

Suppose we have a set of single-particle energy levels $\varepsilon_k$ ($k = 1, 2, \ldots, M$).
The Hamiltonian for $N$ particles is
\begin{equation}
    H = \sum_k \varepsilon_k\, n_k,
\end{equation}
where $n_k$ is the number of particles occupying level $k$.  A microstate is specified by the
set of occupation numbers $\{n_k\}$.

\begin{center}[DIAGRAM: energy level diagram with levels $\varepsilon_1 < \varepsilon_2 < \varepsilon_3 < \ldots$; $n_1$ particles in level 1, $n_2$ in level 2, etc.]\end{center}

The key difference between distinguishable and indistinguishable particles:

\begin{itemize}
    \item If particles are \textbf{indistinguishable}, each set $\{n_k\}$ (with $\sum_k n_k = N$) specifies exactly one microstate.
    \item If particles are \textbf{distinguishable} (each carries a label), the occupation set $\{n_k\}$ corresponds to $N!/\prod_k n_k!$ microstates — the multinomial coefficient counting which labelled particles go in which level.
\end{itemize}

For distinguishable particles, the partition function includes this multiplicity:
\begin{equation}
    Z_N = \sum_{\{n_k\}} \frac{N!}{\prod_k n_k!}\, e^{-\beta\sum_k \varepsilon_k n_k}
        = \sum_{\{n_k\}} \frac{N!}{\prod_k n_k!}\prod_k e^{-\beta\varepsilon_k n_k}.
\end{equation}
Recognising the multinomial expansion $(a_1 + a_2 + \cdots + a_M)^N = \sum_{\{n_k\},\sum n_k=N}\frac{N!}{\prod n_k!}\prod a_k^{n_k}$
with $a_k = e^{-\beta\varepsilon_k}$:
\begin{equation}
    Z_N = \left(e^{-\beta\varepsilon_1} + e^{-\beta\varepsilon_2} + \cdots + e^{-\beta\varepsilon_M}\right)^N
        = \left(\sum_k e^{-\beta\varepsilon_k}\right)^N = (Z_1)^N.
\end{equation}

\begin{formula}[Separable partition function]
\begin{equation}
    Z_N = (Z_1)^N \quad \text{for distinguishable particles when the Hamiltonian is separable.}
\end{equation}
For indistinguishable particles,
\begin{equation}
    Z_N = \frac{1}{N!}(Z_1)^N \quad \text{(valid at high $T$, i.e.\ the classical limit).}
\end{equation}
\end{formula}

The $1/N!$ factor corrects for overcounting of identical-particle microstates.  It is an
approximation (valid when occupation numbers are small, $n_k \ll 1$, so the distinction between
Bose--Einstein and Fermi--Dirac statistics does not matter) and breaks down at low temperatures
where quantum statistics become important.

\chapter{$N$ Independent Harmonic Oscillators}

\section{Setup and Partition Function}

Consider $N$ independent quantum harmonic oscillators.  Each oscillator has energy levels
labeled by a non-negative integer $n_i = 0, 1, 2, \ldots$:
\begin{equation}
    \varepsilon = \hbar\omega\!\left(n_1 + n_2 + n_3 + \tfrac{3}{2}\right),
\end{equation}
where the $\tfrac{3}{2}$ accounts for the three-dimensional zero-point energy.  Since the
particles are distinguishable (they sit at different lattice sites, for example), we have
$Z_N = (Z_1)^N$.  The single-oscillator partition function is
\begin{align}
    Z_1 &= \sum_{n_1,n_2,n_3=0}^{\infty} e^{-\beta\hbar\omega(n_1+n_2+n_3+3/2)} \notag\\
        &= e^{-3\beta\hbar\omega/2} \left(\sum_{n=0}^{\infty} e^{-\beta\hbar\omega n}\right)^3 \notag\\
        &= e^{-3\beta\hbar\omega/2}\left(\frac{1}{1-e^{-\beta\hbar\omega}}\right)^3.
\end{align}
Therefore,
\begin{equation}
    Z_N = (Z_1)^N = \left[\frac{e^{-\frac{1}{2}\beta\hbar\omega}}{1-e^{-\beta\hbar\omega}}\right]^{3N}.
\end{equation}

\section{Internal Energy}

From $U_N = -\partial\log Z_N/\partial\beta$:
\begin{align}
    \log Z_N &= 3N\left[-\tfrac{1}{2}\beta\hbar\omega - \log\!\left(1-e^{-\beta\hbar\omega}\right)\right],\\
    U_N &= -3N\pdv{}{\beta}\!\left[-\tfrac{1}{2}\beta\hbar\omega - \log\!\left(1-e^{-\beta\hbar\omega}\right)\right] \notag\\
        &= 3N\left[\tfrac{1}{2}\hbar\omega + \frac{\hbar\omega\, e^{-\beta\hbar\omega}}{1-e^{-\beta\hbar\omega}}\right].
\end{align}

Rewriting the second term:
\begin{equation}
    \frac{e^{-\beta\hbar\omega}}{1-e^{-\beta\hbar\omega}}
    = \frac{1}{e^{\beta\hbar\omega}-1},
\end{equation}
which is the mean occupation number of a bosonic mode (the Planck factor).  Thus

\begin{formula}[Einstein solid internal energy]
\begin{equation}
    U_N = 3N\hbar\omega\!\left[\frac{1}{2} + \frac{1}{e^{\beta\hbar\omega}-1}\right].
    \label{eq:einstein_U}
\end{equation}
\end{formula}

The first term $\tfrac{3}{2}N\hbar\omega$ is the zero-point energy; the second term gives the
thermal excitation.

\subsection{Checking the Limits}

\paragraph{Low temperature ($T\to 0$):} $e^{\beta\hbar\omega}\gg 1$, so
$1/(e^{\beta\hbar\omega}-1)\approx e^{-\beta\hbar\omega}\to 0$.  Hence $U_N\to\tfrac{3}{2}N\hbar\omega$
(zero-point energy only), as expected.

\paragraph{High temperature ($T\to\infty$):} We need $1/(e^x-1)$ for $x\ll 1$.  Since
$e^x - 1 \approx x + x^2/2 + \cdots$, we have $1/(e^x-1)\approx 1/x - 1/2 + \cdots$.
Thus
\begin{equation}
    U_N \approx 3N\hbar\omega\!\left[\tfrac{1}{2} + \frac{1}{\beta\hbar\omega} - \tfrac{1}{2}\right]
         = 3Nk_BT,
\end{equation}
recovering the classical equipartition result: each of the $3N$ oscillators has kinetic and
potential energy $\tfrac{1}{2}k_BT$ each, totaling $3Nk_BT$.
